% Options for packages loaded elsewhere
\PassOptionsToPackage{unicode}{hyperref}
\PassOptionsToPackage{hyphens}{url}
\documentclass[
]{article}
\usepackage{xcolor}
\usepackage[margin=1in]{geometry}
\usepackage{amsmath,amssymb}
\setcounter{secnumdepth}{-\maxdimen} % remove section numbering
\usepackage{iftex}
\ifPDFTeX
  \usepackage[T1]{fontenc}
  \usepackage[utf8]{inputenc}
  \usepackage{textcomp} % provide euro and other symbols
\else % if luatex or xetex
  \usepackage{unicode-math} % this also loads fontspec
  \defaultfontfeatures{Scale=MatchLowercase}
  \defaultfontfeatures[\rmfamily]{Ligatures=TeX,Scale=1}
\fi
\usepackage{lmodern}
\ifPDFTeX\else
  % xetex/luatex font selection
\fi
% Use upquote if available, for straight quotes in verbatim environments
\IfFileExists{upquote.sty}{\usepackage{upquote}}{}
\IfFileExists{microtype.sty}{% use microtype if available
  \usepackage[]{microtype}
  \UseMicrotypeSet[protrusion]{basicmath} % disable protrusion for tt fonts
}{}
\makeatletter
\@ifundefined{KOMAClassName}{% if non-KOMA class
  \IfFileExists{parskip.sty}{%
    \usepackage{parskip}
  }{% else
    \setlength{\parindent}{0pt}
    \setlength{\parskip}{6pt plus 2pt minus 1pt}}
}{% if KOMA class
  \KOMAoptions{parskip=half}}
\makeatother
\usepackage{color}
\usepackage{fancyvrb}
\newcommand{\VerbBar}{|}
\newcommand{\VERB}{\Verb[commandchars=\\\{\}]}
\DefineVerbatimEnvironment{Highlighting}{Verbatim}{commandchars=\\\{\}}
% Add ',fontsize=\small' for more characters per line
\usepackage{framed}
\definecolor{shadecolor}{RGB}{248,248,248}
\newenvironment{Shaded}{\begin{snugshade}}{\end{snugshade}}
\newcommand{\AlertTok}[1]{\textcolor[rgb]{0.94,0.16,0.16}{#1}}
\newcommand{\AnnotationTok}[1]{\textcolor[rgb]{0.56,0.35,0.01}{\textbf{\textit{#1}}}}
\newcommand{\AttributeTok}[1]{\textcolor[rgb]{0.13,0.29,0.53}{#1}}
\newcommand{\BaseNTok}[1]{\textcolor[rgb]{0.00,0.00,0.81}{#1}}
\newcommand{\BuiltInTok}[1]{#1}
\newcommand{\CharTok}[1]{\textcolor[rgb]{0.31,0.60,0.02}{#1}}
\newcommand{\CommentTok}[1]{\textcolor[rgb]{0.56,0.35,0.01}{\textit{#1}}}
\newcommand{\CommentVarTok}[1]{\textcolor[rgb]{0.56,0.35,0.01}{\textbf{\textit{#1}}}}
\newcommand{\ConstantTok}[1]{\textcolor[rgb]{0.56,0.35,0.01}{#1}}
\newcommand{\ControlFlowTok}[1]{\textcolor[rgb]{0.13,0.29,0.53}{\textbf{#1}}}
\newcommand{\DataTypeTok}[1]{\textcolor[rgb]{0.13,0.29,0.53}{#1}}
\newcommand{\DecValTok}[1]{\textcolor[rgb]{0.00,0.00,0.81}{#1}}
\newcommand{\DocumentationTok}[1]{\textcolor[rgb]{0.56,0.35,0.01}{\textbf{\textit{#1}}}}
\newcommand{\ErrorTok}[1]{\textcolor[rgb]{0.64,0.00,0.00}{\textbf{#1}}}
\newcommand{\ExtensionTok}[1]{#1}
\newcommand{\FloatTok}[1]{\textcolor[rgb]{0.00,0.00,0.81}{#1}}
\newcommand{\FunctionTok}[1]{\textcolor[rgb]{0.13,0.29,0.53}{\textbf{#1}}}
\newcommand{\ImportTok}[1]{#1}
\newcommand{\InformationTok}[1]{\textcolor[rgb]{0.56,0.35,0.01}{\textbf{\textit{#1}}}}
\newcommand{\KeywordTok}[1]{\textcolor[rgb]{0.13,0.29,0.53}{\textbf{#1}}}
\newcommand{\NormalTok}[1]{#1}
\newcommand{\OperatorTok}[1]{\textcolor[rgb]{0.81,0.36,0.00}{\textbf{#1}}}
\newcommand{\OtherTok}[1]{\textcolor[rgb]{0.56,0.35,0.01}{#1}}
\newcommand{\PreprocessorTok}[1]{\textcolor[rgb]{0.56,0.35,0.01}{\textit{#1}}}
\newcommand{\RegionMarkerTok}[1]{#1}
\newcommand{\SpecialCharTok}[1]{\textcolor[rgb]{0.81,0.36,0.00}{\textbf{#1}}}
\newcommand{\SpecialStringTok}[1]{\textcolor[rgb]{0.31,0.60,0.02}{#1}}
\newcommand{\StringTok}[1]{\textcolor[rgb]{0.31,0.60,0.02}{#1}}
\newcommand{\VariableTok}[1]{\textcolor[rgb]{0.00,0.00,0.00}{#1}}
\newcommand{\VerbatimStringTok}[1]{\textcolor[rgb]{0.31,0.60,0.02}{#1}}
\newcommand{\WarningTok}[1]{\textcolor[rgb]{0.56,0.35,0.01}{\textbf{\textit{#1}}}}
\usepackage{graphicx}
\makeatletter
\newsavebox\pandoc@box
\newcommand*\pandocbounded[1]{% scales image to fit in text height/width
  \sbox\pandoc@box{#1}%
  \Gscale@div\@tempa{\textheight}{\dimexpr\ht\pandoc@box+\dp\pandoc@box\relax}%
  \Gscale@div\@tempb{\linewidth}{\wd\pandoc@box}%
  \ifdim\@tempb\p@<\@tempa\p@\let\@tempa\@tempb\fi% select the smaller of both
  \ifdim\@tempa\p@<\p@\scalebox{\@tempa}{\usebox\pandoc@box}%
  \else\usebox{\pandoc@box}%
  \fi%
}
% Set default figure placement to htbp
\def\fps@figure{htbp}
\makeatother
\setlength{\emergencystretch}{3em} % prevent overfull lines
\providecommand{\tightlist}{%
  \setlength{\itemsep}{0pt}\setlength{\parskip}{0pt}}
\usepackage{bookmark}
\IfFileExists{xurl.sty}{\usepackage{xurl}}{} % add URL line breaks if available
\urlstyle{same}
\hypersetup{
  pdftitle={Competing risks and multistate models},
  hidelinks,
  pdfcreator={LaTeX via pandoc}}

\title{Competing risks and multistate models}
\author{}
\date{\vspace{-2.5em}}

\begin{document}
\maketitle

\textbf{Inference lab} using the five-step template: Hypothesis →
Visualise → Assumptions → Conduct → Conclude.

\subsection{Learning objectives}\label{learning-objectives}

\begin{itemize}
\tightlist
\item
  Estimate cause-specific cumulative incidence with \texttt{tidycmprsk}.
\item
  Contrast cause-specific hazards with the Fine-Gray subdistribution
  hazard.
\item
  Set up a three-state illness-death multistate model with
  \texttt{mstate}.
\end{itemize}

\subsection{Prerequisites}\label{prerequisites}

Survival analysis.

\subsection{Background}\label{background}

Competing risks are events that preclude the event of interest: death
from cancer competes with death from other causes. Kaplan- Meier treats
the competing event as censoring, which is wrong when the censoring is
informative. Cumulative incidence functions (CIFs) give the probability
that the event of interest has occurred by time \emph{t} \emph{in the
presence of} competing risks.

Two hazard models are common. Cause-specific Cox models treat each cause
in turn, censoring at other causes; they answer an aetiologic question
(``does the covariate affect the rate of this cause?''). The Fine-Gray
subdistribution hazard model fixes subjects with competing events in the
risk set indefinitely; it answers a prognostic question (``does the
covariate affect the cumulative incidence?'').

Multistate models generalise survival to several possible transitions.
The illness-death model has three states: healthy, ill, dead, with
transitions healthy → ill, healthy → dead, and ill → dead.
\texttt{mstate} builds these objects from tabular data.

Fine-Gray is descriptive; cause-specific Cox is causal (in the simple,
aetiological sense). A covariate can increase the Fine-Gray
subdistribution hazard of cause 1 either by increasing cause 1 itself or
by decreasing cause 2 --- it is a net effect on the cumulative
incidence.

\subsection{Setup}\label{setup}

\begin{Shaded}
\begin{Highlighting}[]
\FunctionTok{library}\NormalTok{(tidyverse)}
\FunctionTok{library}\NormalTok{(survival)}
\FunctionTok{library}\NormalTok{(tidycmprsk)}
\FunctionTok{library}\NormalTok{(ggsurvfit)}
\FunctionTok{library}\NormalTok{(mstate)}
\FunctionTok{set.seed}\NormalTok{(}\DecValTok{42}\NormalTok{)}
\FunctionTok{theme\_set}\NormalTok{(}\FunctionTok{theme\_minimal}\NormalTok{(}\AttributeTok{base\_size =} \DecValTok{12}\NormalTok{))}
\end{Highlighting}
\end{Shaded}

\subsection{1. Hypothesis}\label{hypothesis}

In the \texttt{survival::colon} data (recurrence and death as competing
events), a treatment covariate affects both events. We estimate the
cumulative incidence of each.

\subsection{2. Visualise}\label{visualise}

\begin{Shaded}
\begin{Highlighting}[]
\CommentTok{\# Keep death records (etype == 2) + recurrence merged; for teaching}
\CommentTok{\# we build a two{-}event dataset:}
\NormalTok{dat }\OtherTok{\textless{}{-}}\NormalTok{ colon }\SpecialCharTok{|\textgreater{}}
  \FunctionTok{filter}\NormalTok{(etype }\SpecialCharTok{==} \DecValTok{1}\NormalTok{) }\SpecialCharTok{|\textgreater{}}                       \CommentTok{\# recurrence row per subject}
  \FunctionTok{transmute}\NormalTok{(id, time, }\AttributeTok{status\_rec =}\NormalTok{ status,}
\NormalTok{            rx, age, sex) }\SpecialCharTok{|\textgreater{}}
  \FunctionTok{left\_join}\NormalTok{(}
\NormalTok{    colon }\SpecialCharTok{|\textgreater{}} \FunctionTok{filter}\NormalTok{(etype }\SpecialCharTok{==} \DecValTok{2}\NormalTok{) }\SpecialCharTok{|\textgreater{}}
      \FunctionTok{transmute}\NormalTok{(id, }\AttributeTok{time\_d =}\NormalTok{ time, }\AttributeTok{status\_d =}\NormalTok{ status),}
    \AttributeTok{by =} \StringTok{"id"}
\NormalTok{  ) }\SpecialCharTok{|\textgreater{}}
  \FunctionTok{mutate}\NormalTok{(}
    \CommentTok{\# event type: 0 censored, 1 recurrence, 2 death without recurrence}
    \AttributeTok{etype =} \FunctionTok{case\_when}\NormalTok{(}
\NormalTok{      status\_rec }\SpecialCharTok{==} \DecValTok{1} \SpecialCharTok{\textasciitilde{}} \DecValTok{1}\DataTypeTok{L}\NormalTok{,}
\NormalTok{      status\_d   }\SpecialCharTok{==} \DecValTok{1} \SpecialCharTok{\textasciitilde{}} \DecValTok{2}\DataTypeTok{L}\NormalTok{,}
      \ConstantTok{TRUE}            \SpecialCharTok{\textasciitilde{}} \DecValTok{0}\DataTypeTok{L}
\NormalTok{    ),}
    \AttributeTok{etime =} \FunctionTok{pmin}\NormalTok{(time, time\_d, }\AttributeTok{na.rm =} \ConstantTok{TRUE}\NormalTok{),}
    \AttributeTok{etype\_f =} \FunctionTok{factor}\NormalTok{(etype, }\AttributeTok{levels =} \DecValTok{0}\SpecialCharTok{:}\DecValTok{2}\NormalTok{,}
                     \AttributeTok{labels =} \FunctionTok{c}\NormalTok{(}\StringTok{"censored"}\NormalTok{, }\StringTok{"recurrence"}\NormalTok{, }\StringTok{"death"}\NormalTok{))}
\NormalTok{  ) }\SpecialCharTok{|\textgreater{}}
  \FunctionTok{drop\_na}\NormalTok{(etime, etype\_f, rx)}

\FunctionTok{ggplot}\NormalTok{(dat, }\FunctionTok{aes}\NormalTok{(etime, }\AttributeTok{fill =}\NormalTok{ etype\_f)) }\SpecialCharTok{+}
  \FunctionTok{geom\_histogram}\NormalTok{(}\AttributeTok{bins =} \DecValTok{40}\NormalTok{, }\AttributeTok{alpha =} \FloatTok{0.7}\NormalTok{, }\AttributeTok{position =} \StringTok{"identity"}\NormalTok{) }\SpecialCharTok{+}
  \FunctionTok{labs}\NormalTok{(}\AttributeTok{x =} \StringTok{"Time (days)"}\NormalTok{, }\AttributeTok{y =} \StringTok{"Count"}\NormalTok{, }\AttributeTok{fill =} \ConstantTok{NULL}\NormalTok{)}
\end{Highlighting}
\end{Shaded}

\subsection{3. Assumptions}\label{assumptions}

Non-informative censoring within cause; proportional hazards for
Fine-Gray; for multistate, Markov transitions (future depends only on
current state).

\subsection{4. Conduct}\label{conduct}

\begin{Shaded}
\begin{Highlighting}[]
\NormalTok{cif}
\end{Highlighting}
\end{Shaded}

\begin{Shaded}
\begin{Highlighting}[]
\NormalTok{ggsurvfit}\SpecialCharTok{::}\FunctionTok{ggcuminc}\NormalTok{(cif, }\AttributeTok{outcome =} \FunctionTok{c}\NormalTok{(}\StringTok{"recurrence"}\NormalTok{, }\StringTok{"death"}\NormalTok{)) }\SpecialCharTok{+}
  \FunctionTok{labs}\NormalTok{(}\AttributeTok{x =} \StringTok{"Days"}\NormalTok{, }\AttributeTok{y =} \StringTok{"Cumulative incidence"}\NormalTok{)}
\end{Highlighting}
\end{Shaded}

\begin{Shaded}
\begin{Highlighting}[]
\NormalTok{fg }\OtherTok{\textless{}{-}} \FunctionTok{crr}\NormalTok{(}\FunctionTok{Surv}\NormalTok{(etime, etype\_f) }\SpecialCharTok{\textasciitilde{}}\NormalTok{ rx }\SpecialCharTok{+}\NormalTok{ age }\SpecialCharTok{+}\NormalTok{ sex, }\AttributeTok{data =}\NormalTok{ dat,}
          \AttributeTok{failcode =} \StringTok{"recurrence"}\NormalTok{)}
\NormalTok{fg}
\end{Highlighting}
\end{Shaded}

\begin{Shaded}
\begin{Highlighting}[]
\NormalTok{n }\OtherTok{\textless{}{-}} \DecValTok{300}
\NormalTok{sim }\OtherTok{\textless{}{-}} \FunctionTok{tibble}\NormalTok{(}
  \AttributeTok{id =} \FunctionTok{seq\_len}\NormalTok{(n),}
  \AttributeTok{t\_ill   =} \FunctionTok{rexp}\NormalTok{(n, }\AttributeTok{rate =} \FloatTok{0.02}\NormalTok{),}
  \AttributeTok{t\_death =} \FunctionTok{rexp}\NormalTok{(n, }\AttributeTok{rate =} \FloatTok{0.01}\NormalTok{)}
\NormalTok{) }\SpecialCharTok{|\textgreater{}}
  \FunctionTok{mutate}\NormalTok{(}
    \AttributeTok{t1 =} \FunctionTok{pmin}\NormalTok{(t\_ill, t\_death, }\DecValTok{50}\NormalTok{),}
    \AttributeTok{s1 =} \FunctionTok{if\_else}\NormalTok{(t\_ill }\SpecialCharTok{\textless{}}\NormalTok{ t\_death }\SpecialCharTok{\&}\NormalTok{ t\_ill }\SpecialCharTok{\textless{}=} \DecValTok{50}\NormalTok{, }\DecValTok{1}\DataTypeTok{L}\NormalTok{, }\DecValTok{0}\DataTypeTok{L}\NormalTok{),}
    \AttributeTok{s2 =} \FunctionTok{if\_else}\NormalTok{(t\_death }\SpecialCharTok{\textless{}=} \DecValTok{50} \SpecialCharTok{\&}\NormalTok{ t\_death }\SpecialCharTok{\textless{}}\NormalTok{ t\_ill, }\DecValTok{1}\DataTypeTok{L}\NormalTok{, }\DecValTok{0}\DataTypeTok{L}\NormalTok{)}
\NormalTok{  )}

\NormalTok{tmat }\OtherTok{\textless{}{-}} \FunctionTok{transMat}\NormalTok{(}\AttributeTok{x =} \FunctionTok{list}\NormalTok{(}\FunctionTok{c}\NormalTok{(}\DecValTok{2}\NormalTok{, }\DecValTok{3}\NormalTok{), }\FunctionTok{c}\NormalTok{(}\DecValTok{3}\NormalTok{), }\FunctionTok{c}\NormalTok{()),}
                 \AttributeTok{names =} \FunctionTok{c}\NormalTok{(}\StringTok{"healthy"}\NormalTok{, }\StringTok{"ill"}\NormalTok{, }\StringTok{"dead"}\NormalTok{))}
\NormalTok{tmat}
\end{Highlighting}
\end{Shaded}

\subsection{5. Concluding statement}\label{concluding-statement}

\begin{quote}
Cause-specific cumulative incidence functions differed by treatment arm
in the colon data, with the Fine-Gray subdistribution hazard estimated
above. The three-state illness- death transition matrix (healthy → ill →
dead) shows how \texttt{mstate} represents the problem; the full fit is
straightforward once the long-format dataset is built via
\texttt{msprep}.
\end{quote}

\subsection{Common pitfalls}\label{common-pitfalls}

\begin{itemize}
\tightlist
\item
  Reporting 1 − KM as the cumulative incidence in the presence of
  competing risks.
\item
  Using Fine-Gray when the question is aetiological.
\item
  Reporting both cause-specific and Fine-Gray hazard ratios without
  distinguishing them.
\item
  Assuming Markov transitions in a multistate model when history
  matters.
\end{itemize}

\subsection{Further reading}\label{further-reading}

\begin{itemize}
\tightlist
\item
  Fine JP, Gray RJ (1999), \emph{A proportional hazards model for the
  subdistribution of a competing risk}.
\item
  Putter H, Fiocco M, Geskus RB (2007), \emph{Tutorial in biostatistics:
  competing risks and multi-state models}.
\item
  Austin PC, Fine JP (2017), \emph{Practical recommendations for
  reporting Fine-Gray model analyses}.
\end{itemize}

\subsection{Session info}\label{session-info}

\begin{Shaded}
\begin{Highlighting}[]

\end{Highlighting}
\end{Shaded}

\subsection{See also --- chapter index}\label{see-also-chapter-index}

\begin{itemize}
\tightlist
\item
  \href{../index.Rmd}{Methods in this chapter}
\item
  \href{../../../ch00-find-your-method.Rmd}{Find your method (Chapter
  0)}
\end{itemize}

\end{document}
